%---------------------------------------------------------------------
\chapter{Requirements and System Analysis}
%---------------------------------------------------------------------

This chapter translates the problem stated in Chapter~1 into concrete requirements. Each functional
requirement is numbered so that it can be traced to the design in Chapter~4 and tested in Chapter~7,
and each non-functional requirement is given a numerical bound.

\section{Stakeholders and Users}

Four groups interact with or are affected by LoomTrace.

\begin{itemize}[itemsep=2pt]
\item \textbf{Requesters (primary users).} Buyers of handloom products, marketplace moderators who
review seller listings, and retail staff who receive stock described as GI-tagged. They need a quick,
understandable answer with the reasons and sources behind it, and a reference number with which to
follow up.
\item \textbf{Reviewers.} Staff of a weavers' cooperative society, a handloom development office or a
marketplace trust-and-safety team who resolve cases that the agent cannot decide. They need the
extracted claim, the reason for escalation and a way to record a decision with a note.
\item \textbf{Weavers and GI proprietors (affected parties).} Producers within the registered area
benefit when false claims are detected, but are harmed if a genuine product is wrongly marked as
conflicting. Their interest is the main reason that the system abstains rather than guesses.
\item \textbf{Course evaluators and developers.} They need step-level traces, reproducible evaluation
scripts and clear documentation of what is real and what is simulated.
\end{itemize}

\section{Functional Requirements}

\begin{description}[leftmargin=1.6cm,style=nextline,itemsep=2pt]
\item[FR1] The system shall accept a product claim as free label text containing any subset of the
fields product, region, state, material, weaver identifier and batch identifier, and optionally a
label image.
\item[FR2] The system shall extract these fields into a typed structure and resolve local product
names and aliases to the registered GI name.
\item[FR3] The system shall decide which checks are applicable to the claim and run them using tools
over the GI registry, GI document text, census data, weaver records and batch records.
\item[FR4] The system shall output exactly one of three verdicts, \emph{Supported}, \emph{Conflicting
evidence} or \emph{Insufficient evidence}, together with the result of every check and the source
identifiers that support it.
\item[FR5] The system shall re-plan and retrieve again when the product cannot be resolved or the
retrieved GI text is weak, for at most three iterations in total.
\item[FR6] The system shall send every \emph{Insufficient evidence} verdict, and any verdict without a
citation or below a confidence threshold, to a human review queue, and shall allow a requester to ask
for review of any verdict.
\item[FR7] The system shall allow a reviewer to record a final decision with a note, and shall allow
the requester to see that decision using the case reference.
\item[FR8] The system shall store the full case state after every node, and keep a history of all
cases and decisions.
\item[FR9] The system shall report claims for products outside the covered GIs as ``not covered''
rather than as conflicting.
\end{description}

\section{Non-Functional Requirements}

\begin{description}[leftmargin=1.6cm,style=nextline,itemsep=2pt]
\item[NFR1] \textbf{Latency.} A verdict shall be returned in under 15 seconds per case, the target set
in the project's design document.
\item[NFR2] \textbf{Bounded cost.} The agent shall execute at most three planning iterations per case,
and therefore at most fifteen tool calls; in rule mode it shall make no paid API calls.
\item[NFR3] \textbf{Traceability.} Every node execution shall be recorded with its inputs and outputs,
and every verdict shall list its source identifiers.
\item[NFR4] \textbf{Safety.} No \emph{Supported} or \emph{Conflicting} verdict shall be issued without at
least one citation and a grade of at least 0.6; otherwise the agent shall abstain.
\item[NFR5] \textbf{Quality.} Verdict accuracy on the labelled test set shall be at least 80\%, and at
least 90\% of cases labelled \emph{Insufficient evidence} shall be routed to human review.
\item[NFR6] \textbf{Reproducibility.} All reported numbers shall be reproducible by running the
scripts in the repository on a laptop without internet access.
\end{description}

\section{Data Sources}

Table~\ref{tab:data} lists the data sources used. Two of them are official Government of India
publications, one is the IP India registration information recorded in the project's design document,
and the producer-level records are simulated because individual weaver and batch data are not
published.

\begin{table}[H]
\centering
\caption{Data sources used in LoomTrace}
\label{tab:data}
\small
\begin{tabularx}{\linewidth}{@{}>{\raggedright\arraybackslash}p{3.0cm} >{\raggedright\arraybackslash}X >{\raggedright\arraybackslash}p{2.1cm} >{\raggedright\arraybackslash}p{3.1cm}@{}}
\toprule
\textbf{Source} & \textbf{Description and use} & \textbf{Size} & \textbf{Origin} \\ \midrule
GI registry entries (3 products) & GI application number, state, districts and materials for Kancheepuram Silk (GI~15), Pochampalli Ikat (GI~4) and Banaras Brocades and Sarees (GI~99); registry checks & 3 records & IP India registration list~\cite{ipindia}, as recorded in the design document \\
GI catalogue (12 products) & Place of production, materials and description of 12 further handloom GIs; registry checks and retrieval corpus & 12 records, 12 passages & PIB / TRIFED, Nov.~2021~\cite{pibgi} \\
GI description passages (3 products) & Short summaries of the registered area, method and materials; retrieval corpus & 9 passages & Written by the team from the GI descriptions; page numbers indicative \\
Handloom census & Handloom weavers and allied workers per state/UT; plausibility check & 31 rows & PIB annexure, release ID 2202869~\cite{pibcensus} \\
Weaver registry & Weaver identifier, name, cluster, district, looms & 180 records & \textbf{Simulated} (seeded generator) \\
Production batches & Batch identifier, weaver, loom, yarn, start and finish dates & 360 records & \textbf{Simulated} (seeded generator) \\
Golden test set & Labelled claims, 3 products $\times$ 3 verdicts $\times$ 5 & 45 cases & Constructed by the team \\
Demonstration tests & Hand-written claims covering 9 products & 15 cases & Constructed by the team \\
\bottomrule
\end{tabularx}
\end{table}

The government data were obtained by downloading the two PIB PDF documents and extracting their text.
The census table was transcribed in full, including all 31 states and union territories. From the GI
catalogue, only handloom textile entries whose place of production is stated in the document were
used; two handloom entries, Khandua Saree and Karvath Kati Sarees, were excluded because the catalogue
does not name their area of production. The catalogue does not print GI application numbers, so for
these twelve products the system reports ``Registered GI'' without a number rather than supplying a
number from an unverified source.

The weaver and batch records are generated by a seeded pseudo-random generator (seed 42), so the same
records are produced on every run. For each product, 12 weavers are assigned in rotation to the
districts of that product's GI area, and each weaver is given two batches whose yarn is drawn from the
product's registered materials. These records are labelled as simulated in the source code, in the
tool outputs (``SIMULATED record'') and in the user interface.

\section{Hardware and Software Requirements}

Development and all evaluation runs used an Apple MacBook with an M2 processor and 8~GB of memory,
running macOS 15.6. The software is written in Python~3.13 and uses LangGraph for orchestration,
Pydantic for typed tool schemas, SQLite for persistence, Streamlit for the interface and pandas for
tabular display. The optional language-model path uses the Groq API with the Llama~3.1~8B Instant
model, and the optional OCR path uses EasyOCR. Exact versions are given in Table~\ref{tab:stack} and
Appendix~\ref{app:run}. No graphics processor is required.

\section{Feasibility}

\paragraph{Technical feasibility.} All checks reduce to lookups and comparisons over small tables and a
corpus of 21 passages, so they can be executed in milliseconds on a laptop. The only components that
require external services, the language model and OCR, are optional and are isolated behind tool
interfaces with deterministic fall-backs.

\paragraph{Cost.} All results in this report were produced in rule mode, without an API key, and
therefore made zero language-model calls and incurred no cost. In the optional model mode, the design
limits model usage to two short calls per case (field extraction and verdict explanation); the model
selected is available on the free tier of the Groq service~\cite{groq}.

\paragraph{Time and scope control.} The scope was limited to verification against stored reference data;
live access to government portals was avoided because the IP India application-detail pages returned
errors during data collection and because the evaluation must be reproducible offline. The project
was initially scoped to three products, as specified in the design document, and was extended to
fifteen products once the official GI catalogue had been obtained.
